
In Linux, a \cc{dev_t} type is used to represent device numbers. It is \cc{typedef}ed to
a 32-bit unsigned integer (\cc{u32}). It is composed of two parts: a major number and a minor
number. Within the 32 bits, 12 bits are set aside for the major number, and the remaining
20 bits for the minor number. In C, to construct a \cc{dev_t} or to fetch the device
numbers, users need to use a some macros, which are implemented with convoluted bit twiddling:
\begin{lstlisting}[language=C]
#define MINORBITS    20
#define MINORMASK    ((1U << MINORBITS) - 1)
#define MAJOR(dev)   ((unsigned int) ((dev) >> MINORBITS)) // get major
#define MINOR(dev)   ((unsigned int) ((dev) & MINORMASK))  // get minor
#define MKDEV(ma,mi) (((ma) << MINORBITS) | (mi))          // construct dev_t
\end{lstlisting}
Logically, the \cc{dev_t} type is simply a product of two integers.
To implement it in a high-level language, users would typically define it as a record
type with two integer fields. For example, in \cogent, it may be defined as:\footnote{
As an opening example, we show pseudo-\cogent code here to omit the irrelevant \cogent
intricacies.}

\begin{lstlisting}[language=Cogent]
type Dev_t = { major : Int, minor : Int }
\end{lstlisting}
and to access the fields of such an object \cg{dev},
the normal member operation (\cg{dev.major} and \cg{dev.minor}) can be used.
With this abtract view of the type, which facilitates simpler operations on this type,
the user unfortunately has to hand over control of the type's layout to the
\cogent compiler. The user then needs to study the compiler's code generation convention
and write data format conversion functions to related the compiled \cg{Dev_t} type with
the native C type \cc{dev_t}. In fact, the user does not have to make this compromise. With
\dargent, the user can have both the abstraction and control over types' memory
layouts. By specifying a \dargent layout description, the \cogent compiler can
generate C code with the specified layout:
\begin{lstlisting}[language=Cogent]
type Dev_t = { major : Int, minor : Int }
  layout record { major : 12b, minor : 20b }
\end{lstlisting}
This is transparent to the functional smenatics
of the \cogent program. Namely, the \cg{dev.major} operation continues to work regardless
of the underlying layout of \cg{dev}.

